% =================================================================
%  CSE 815 -- Machine Learning
%  Module 2, Week 4, Part 1
%  Neural Networks: Neurons, Layers, and Forward Propagation
% =================================================================
\documentclass[aspectratio=169,10pt]{beamer}
\usepackage{cse815}

\title[CSE 815 -- Week 4, Part 1]{Neural Networks}
\subtitle{Neurons, Layers, and Making Predictions by Forward Propagation}
\author{Rokan Uddin Faruqui\\ \small Professor, Department of CSE}
\institute{\small University of Chittagong \\ \texttt{rufaruqui@cu.ac.bd}}
\date[Week 4 / Part 1]{Week 4 $\cdot$ Part 1}

\begin{document}

\begin{frame}[plain,noframenumbering]
  \titlepage
\end{frame}

% ---------------------------------------------------------------
\begin{frame}{From one logistic unit to many}
  \begin{columns}[T,onlytextwidth]
    \begin{column}{0.5\textwidth}
      \begin{block}{Module 1 gave us}
        \[ f_{\vw,b}(\vx)=g(\vw\cdot\vx+b),\qquad g(z)=\frac{1}{1+e^{-z}} \]
        a classifier whose decision boundary is a line --- unless we hand-craft features such
        as $x_1^2$, $x_1x_2$ to bend it.
      \end{block}
      \footnotesize
      Hand-crafting features does not scale to an image with a million pixels. Module 2 asks:
      can the model \emph{learn} its own features? A neural network does exactly that, and
      every one of its units is the logistic regression you already know.
    \end{column}
    \begin{column}{0.46\textwidth}
      \begin{keyidea}[Today]
        \footnotesize
        \begin{enumerate}
          \item Neurons and the brain --- and where the analogy stops
          \item Demand prediction: a network as stacked logistic units
          \item Recognising images: features learned layer by layer
          \item Inside a layer: notation and computation
          \item Deeper networks and forward propagation
        \end{enumerate}
      \end{keyidea}
      \footnotesize Part 2 turns all of this into code: TensorFlow, then NumPy from scratch,
      then matrix multiplication.
    \end{column}
  \end{columns}
\end{frame}

% ===============================================================
\section{Neurons and the Brain}

\begin{frame}{A short history}
  \footnotesize
  \begin{center}
  \begin{tikzpicture}[x=1.4cm]
    \draw[cugray,thick,-{Stealth}] (0,0) -- (10.3,0);
    \foreach \x/\yr in {0.3/1943,1.6/1958,3.4/1969,5.2/1986,7.0/1998,8.8/2012}{
      \fill[cuteal] (\x,0) circle (2pt);
      \node[below=2pt,font=\scriptsize\bfseries,cublue] at (\x,0) {\yr};
    }
    \node[above,align=center,font=\scriptsize,text width=1.6cm] at (0.3,0.1) {McCulloch \& Pitts: a logical neuron};
    \node[above,align=center,font=\scriptsize,text width=1.7cm] at (1.6,0.1) {Rosenblatt's perceptron};
    \node[above,align=center,font=\scriptsize,text width=1.7cm] at (3.4,0.1) {Minsky \& Papert: limits of one layer};
    \node[above,align=center,font=\scriptsize,text width=1.8cm] at (5.2,0.1) {Back-propagation popularised};
    \node[above,align=center,font=\scriptsize,text width=1.7cm] at (7.0,0.1) {LeNet reads cheques};
    \node[above,align=center,font=\scriptsize,text width=1.8cm] at (8.8,0.1) {AlexNet wins ImageNet};
  \end{tikzpicture}
  \end{center}
  \vspace{1mm}
  \begin{columns}[T,onlytextwidth]
    \begin{column}{0.5\textwidth}
      \begin{block}{Original motivation}
        Write software that mimics how the brain learns. Interest rose, fell, and rose again
        --- twice.
      \end{block}
      \begin{block}{Today's motivation}
        Engineering, not biology: networks are simply the most accurate known models for
        speech, images and text, once there is enough data.
      \end{block}
    \end{column}
    \begin{column}{0.46\textwidth}
      \begin{block}{Where they won first}
        \begin{enumerate}
          \item Speech recognition
          \item Computer vision
          \item Natural-language processing
        \end{enumerate}
        and now climate, medicine, advertising, recommendation \dots
      \end{block}
    \end{column}
  \end{columns}
\end{frame}

\begin{frame}{The biological neuron, and the artificial one}
  \begin{columns}[T,onlytextwidth]
    \begin{column}{0.5\textwidth}
      \centering {\scriptsize Biological: dendrites $\to$ cell body $\to$ axon}\par
      \begin{tikzpicture}[scale=0.85]
        \fill[cuteal!15,draw=cuteal,thick] (0,0) circle (0.55);
        \node[font=\tiny,cuteal] at (0,0) {nucleus};
        \foreach \a in {120,150,180,210,240}{
          \draw[cuteal,thick] (\a:0.55) -- (\a:1.6);
          \draw[cuteal] (\a:1.2) -- ++(\a+25:0.35);
        }
        \draw[cuteal,very thick] (0.55,0) -- (3.0,0);
        \foreach \a in {-30,0,30}{\draw[cuteal,thick] (3.0,0) -- ++(\a:0.5);}
        \node[font=\tiny,cugray] at (-1.5,1.25) {dendrites (inputs)};
        \node[font=\tiny,cugray] at (1.8,0.25) {axon (output)};
      \end{tikzpicture}
      \vspace{1mm}

      \centering {\scriptsize Artificial: inputs $\to$ weighted sum $\to$ activation}\par
      \begin{tikzpicture}[node distance=5mm]
        \node[inneuron] (x1) at (0,0.6) {$x_1$};
        \node[inneuron] (x2) at (0,0) {$x_2$};
        \node[inneuron] (x3) at (0,-0.6) {$x_3$};
        \node[neuron,minimum size=9mm] (n) at (2.2,0) {$g(z)$};
        \node[font=\scriptsize] (a) at (3.8,0) {$a$};
        \foreach \i in {1,2,3}{\draw[synapse] (x\i) -- node[pos=0.35,above,font=\tiny,sloped]{$w_\i$} (n);}
        \draw[synapse] (n) -- (a);
      \end{tikzpicture}
    \end{column}
    \begin{column}{0.46\textwidth}
      \begin{defnbox}[Artificial neuron (unit)]
        \small
        \[ z=\vw\cdot\vx+b,\qquad a=g(z) \]
        $a$ is the unit's \textbf{activation} --- the number it sends on. With the sigmoid
        $g$, one unit \emph{is} logistic regression.
      \end{defnbox}
      \begin{pitfall}[A loose analogy]
        \small
        We still do not know how the brain learns. Modern networks owe their success to
        mathematics and data, not to fidelity to neurobiology. Do not argue from biology
        when designing a model.
      \end{pitfall}
    \end{column}
  \end{columns}
\end{frame}

\begin{frame}{Why now? Scale}
  \begin{columns}[T,onlytextwidth]
    \begin{column}{0.52\textwidth}
      \centering {\scriptsize Schematic --- not measured data}\par
      \begin{tikzpicture}
        \begin{axis}[width=7.2cm,height=4.8cm,xlabel={amount of data},ylabel={performance},
          xmin=0,xmax=10,ymin=0,ymax=10,xtick=\empty,ytick=\empty,
          label style={font=\scriptsize},
          legend style={font=\tiny,at={(0.98,0.04)},anchor=south east,draw=none,fill=white},
          legend cell align=left]
          \addplot[cugray,very thick,domain=0:10,samples=60]{3.3*(1-exp(-0.9*x))};
          \addlegendentry{logistic / linear regression}
          \addplot[cuteal,very thick,domain=0:10,samples=60]{5.8*(1-exp(-0.55*x))};
          \addlegendentry{small network}
          \addplot[cuamber,very thick,domain=0:10,samples=60]{7.4*(1-exp(-0.38*x))};
          \addlegendentry{medium network}
          \addplot[curose,very thick,domain=0:10,samples=60]{9.4*(1-exp(-0.26*x))};
          \addlegendentry{large network}
        \end{axis}
      \end{tikzpicture}
    \end{column}
    \begin{column}{0.44\textwidth}
      \footnotesize
      \begin{block}{Two things changed}
        \begin{itemize}
          \item \textbf{Data.} Digitised records, phones and the web made labelled data
                abundant.
          \item \textbf{Compute.} GPUs made the matrix arithmetic of large networks cheap
                --- the subject of Part 2.
        \end{itemize}
      \end{block}
      \begin{keyidea}
        Classical models plateau as data grows. Larger networks keep improving --- so, given
        the data, a bigger network is often the most reliable route to better accuracy.
      \end{keyidea}
    \end{column}
  \end{columns}
\end{frame}

% ===============================================================
\section{Demand Prediction}

\begin{frame}{One neuron: will this T-shirt be a top seller?}
  \begin{columns}[T,onlytextwidth]
    \begin{column}{0.5\textwidth}
      \centering {\scriptsize $x$ = price; $y=1$ if a top seller}\par
      \begin{tikzpicture}
        \begin{axis}[width=6.2cm,height=3.9cm,xlabel={price (Tk)},ylabel={$a$},
          xmin=0,xmax=1000,ymin=-0.1,ymax=1.1,ytick={0,0.5,1},
          tick label style={font=\scriptsize},label style={font=\scriptsize}]
          \addplot[only marks,mark=*,mark size=1.6pt,cublue,forget plot]
            coordinates {(150,1)(220,1)(300,1)(380,1)(450,1)(560,0)(640,0)(720,0)(850,0)(930,0)};
          \addplot[curose,very thick,domain=0:1000,samples=80,forget plot]
            {1/(1+exp(0.02*(x-500)))};
        \end{axis}
      \end{tikzpicture}
    \end{column}
    \begin{column}{0.46\textwidth}
      \[ a = f(x) = \frac{1}{1+e^{-(wx+b)}} \]
      \footnotesize
      \begin{center}
      \begin{tikzpicture}
        \node[inneuron] (x) at (0,0) {$x$};
        \node[neuron,minimum size=8mm] (n) at (2,0) {};
        \node[font=\scriptsize] (a) at (3.5,0) {$a$};
        \draw[synapse] (x) -- node[above,font=\tiny]{price} (n);
        \draw[synapse] (n) -- node[above,font=\tiny]{prob.} (a);
      \end{tikzpicture}
      \end{center}
      \begin{block}{Vocabulary}
        The circle is a \textbf{neuron}. Its output $a$ is its \textbf{activation} --- here
        the estimated probability of being a top seller. The term comes from biology: how
        strongly a neuron ``fires''.
      \end{block}
    \end{column}
  \end{columns}
\end{frame}

\begin{frame}{Four features, three intermediate ideas, one answer}
  \begin{columns}[T,onlytextwidth]
    \begin{column}{0.56\textwidth}
      \centering
      \begin{tikzpicture}[x=2.1cm,y=1.05cm]
        \foreach \n/\lab [count=\i] in {price,shipping cost,marketing,material}{
          \node[inneuron] (x\i) at (0,-\i) {};
          \node[left,font=\tiny] at (x\i.west) {\lab};
        }
        \foreach \n [count=\j] in {affordability,awareness,perceived quality}{
          \node[neuron] (h\j) at (1,-\j-0.5) {};
          \node[above=1pt,font=\tiny,cublue,fill=white,inner sep=1pt] at (h\j.north) {\n};
        }
        \node[outneuron] (o) at (2,-2.5) {};
        \node[right,font=\tiny] (p) at (2.3,-2.5) {$P(\text{top seller})$};
        \draw[synapse] (o) -- (p);
        \foreach \i in {1,...,4}{\foreach \j in {1,2,3}{\draw[synapse] (x\i) -- (h\j);}}
        \foreach \j in {1,2,3}{\draw[synapse] (h\j) -- (o);}
        \node[font=\tiny,cugray] at (0,-5) {input layer $\vx$};
        \node[font=\tiny,cugray] at (1,-5) {hidden layer $\va$};
        \node[font=\tiny,cugray] at (2,-5) {output layer};
      \end{tikzpicture}
    \end{column}
    \begin{column}{0.42\textwidth}
      \footnotesize
      \begin{block}{How to read it}
        We \emph{imagine} that the hidden units estimate affordability (from price and
        shipping), awareness (from marketing), and perceived quality (from price and
        material). The output unit is a logistic regression on those three numbers.
      \end{block}
      \begin{keyidea}[Full connectivity]
        In practice every hidden unit sees \emph{every} input. We do not wire affordability
        to price by hand --- training decides which inputs each unit should use, by setting
        the weights.
      \end{keyidea}
    \end{column}
  \end{columns}
\end{frame}

\begin{frame}{The network learns its own features}
  \begin{columns}[T,onlytextwidth]
    \begin{column}{0.5\textwidth}
      \begin{block}{Week 2: feature engineering}
        \footnotesize
        You decided the useful feature was $\text{area}=\text{frontage}\times\text{depth}$,
        computed it, and handed it to linear regression.
      \end{block}
      \begin{block}{Week 4: feature learning}
        \footnotesize
        The hidden layer computes features $\va$ from $\vx$; the output layer is a logistic
        regression on $\va$. What the features \emph{are} is chosen by gradient descent, not
        by you.
      \end{block}
      \begin{center}
        \footnotesize
        \begin{tikzpicture}[node distance=4mm]
          \node[draw=cuteal,fill=cuteal!12,rounded corners,inner sep=3pt] (x) {$\vx$};
          \node[draw=cublue,fill=culight,rounded corners,inner sep=3pt,right=8mm of x] (a) {$\va$};
          \node[draw=cuamber,fill=cuamber!15,rounded corners,inner sep=3pt,right=8mm of a] (y) {$a\approx P(y=1)$};
          \draw[synapse] (x) -- node[above,font=\tiny]{layer 1} (a);
          \draw[synapse] (a) -- node[above,font=\tiny]{layer 2} (y);
        \end{tikzpicture}
      \end{center}
    \end{column}
    \begin{column}{0.46\textwidth}
      \begin{defnbox}[Architecture]
        \small
        The choices you make \emph{before} training: how many hidden layers, and how many
        units in each. Weights and biases are what training then learns.
      \end{defnbox}
      \footnotesize
      \textbf{Terminology.}
      \begin{itemize}
        \item \textbf{Input layer}: the feature vector $\vx$ --- it computes nothing.
        \item \textbf{Hidden layer}: its correct values never appear in the training set,
              which contains only $(\vx,y)$.
        \item \textbf{Output layer}: produces the prediction.
        \item A network with several hidden layers is a \textbf{multilayer perceptron}.
      \end{itemize}
    \end{column}
  \end{columns}
\end{frame}

% ===============================================================
\section{Example: Recognising Images}

\begin{frame}{An image is a long vector}
  \begin{columns}[T,onlytextwidth]
    \begin{column}{0.4\textwidth}
      \centering
      \begin{tikzpicture}[scale=0.32]
        \foreach \i in {0,...,7}{\foreach \j in {0,...,7}{
          \pgfmathsetmacro{\v}{mod(37*\i+11*\j+5*\i*\j,90)}
          \fill[cublue!\v] (\i,\j) rectangle ++(1,1);
        }}
        \draw[cugray] (0,0) grid (8,8);
        \node[font=\tiny,cugray] at (4,-0.8) {each cell: brightness $0$--$255$};
      \end{tikzpicture}
      \vspace{1mm}

      \footnotesize
      A $1000\times1000$ greyscale photograph is a vector
      \[ \vx\in\R^{1\,000\,000}. \]
      A colour photo is three times that.
    \end{column}
    \begin{column}{0.56\textwidth}
      \begin{block}{Face recognition}
        \footnotesize
        Input: the pixel vector $\vx$. Output: the identity of the person. Logistic
        regression on raw pixels fails --- no single pixel says anything about identity, and
        no human can engineer a million-pixel feature set.
      \end{block}
      \begin{keyidea}[What a trained network turns out to compute]
        \footnotesize
        Inspecting a trained face network, researchers find:
        \begin{enumerate}
          \item first hidden layer: short \textbf{edges} at various angles;
          \item second: \textbf{parts} --- an eye, the corner of a nose;
          \item third: coarse \textbf{whole-face} shapes.
        \end{enumerate}
        Nobody told it to look for edges. Each layer builds on the one before.
      \end{keyidea}
    \end{column}
  \end{columns}
\end{frame}

\begin{frame}{Increasingly abstract features, learned from data}
  \centering
  \begin{tikzpicture}[
      stage/.style={draw=cublue,fill=culight,rounded corners,minimum width=2.3cm,
                    minimum height=1.6cm,align=center,font=\scriptsize}]
    \node[stage,draw=cuteal,fill=cuteal!10] (in) at (0,0) {pixels\\[2pt]$\vx$};
    \node[stage] (l1) at (3.2,0) {edges\\[2pt]small windows};
    \node[stage] (l2) at (6.4,0) {parts\\[2pt]larger windows};
    \node[stage] (l3) at (9.6,0) {objects\\[2pt]whole image};
    \node[stage,draw=cuamber,fill=cuamber!12] (out) at (12.6,0) {answer\\[2pt]$P(\text{person }k)$};
    \foreach \a/\b in {in/l1,l1/l2,l2/l3,l3/out}{\draw[synapse,thick] (\a) -- (\b);}
    % edge glyphs
    \foreach \ang [count=\k] in {0,45,90,135}{
      \draw[cublue,thick] ($(l1.south)+(-0.9+0.6*\k-0.6,-0.45)$) -- ++(\ang:0.35);
    }
    \draw[cublue,thick] ($(l2.south)+(-0.5,-0.55)$) circle (0.15);
    \draw[cublue,thick] ($(l2.south)+(0.3,-0.55)$) arc (180:360:0.25 and 0.12);
    \draw[cublue,thick] ($(l3.south)+(0,-0.6)$) ellipse (0.3 and 0.38);
  \end{tikzpicture}
  \vspace{4mm}
  \begin{columns}[T,onlytextwidth]
    \begin{column}{0.5\textwidth}
      \begin{block}{Same architecture, different data}
        \footnotesize
        Train the identical network on photographs of cars and the hidden layers learn
        wheels and windows instead of eyes and noses. The architecture is generic; the
        features come from the data.
      \end{block}
    \end{column}
    \begin{column}{0.46\textwidth}
      \begin{pitfall}[Interpretations are post hoc]
        \footnotesize
        ``Layer 2 detects eyes'' is a description we attach after training, by visualising
        what excites each unit. It is not a guarantee, and many units have no tidy label.
      \end{pitfall}
    \end{column}
  \end{columns}
\end{frame}

% ===============================================================
\section{Inside a Neural Network Layer}

\begin{frame}{Layer 1, unit by unit}
  \begin{columns}[T,onlytextwidth]
    \begin{column}{0.4\textwidth}
      \centering
      \begin{tikzpicture}[x=1.9cm,y=0.95cm]
        \node[inneuron] (x1) at (0,-1) {$x_1$};
        \node[inneuron] (x2) at (0,-2) {$x_2$};
        \foreach \j in {1,2,3}{\node[neuron,minimum size=8mm] (h\j) at (1,-\j+0.5) {\tiny $a^{[1]}_{\j}$};}
        \node[outneuron,minimum size=8mm] (o) at (2,-1.5) {\tiny $a^{[2]}$};
        \foreach \i in {1,2}{\foreach \j in {1,2,3}{\draw[synapse] (x\i) -- (h\j);}}
        \foreach \j in {1,2,3}{\draw[synapse] (h\j) -- (o);}
        \node[font=\tiny,cugray] at (0,-3.4) {layer 0};
        \node[font=\tiny,cugray] at (1,-3.4) {layer 1};
        \node[font=\tiny,cugray] at (2,-3.4) {layer 2};
      \end{tikzpicture}
      \vspace{1mm}

      \footnotesize
      Each hidden unit $j$ has its own parameters $\vw^{[1]}_j$, $b^{[1]}_j$, and is a
      logistic regression on $\vx$.
    \end{column}
    \begin{column}{0.56\textwidth}
      \small
      \begin{align*}
        a^{[1]}_1 &= g\bigl(\vw^{[1]}_1\cdot\vx + b^{[1]}_1\bigr)\\
        a^{[1]}_2 &= g\bigl(\vw^{[1]}_2\cdot\vx + b^{[1]}_2\bigr)\\
        a^{[1]}_3 &= g\bigl(\vw^{[1]}_3\cdot\vx + b^{[1]}_3\bigr)
      \end{align*}
      \vspace{-3mm}
      \[ \va^{[1]} = \bigl[a^{[1]}_1,\ a^{[1]}_2,\ a^{[1]}_3\bigr] \]
      \footnotesize
      The output layer then does the same thing to $\va^{[1]}$:
      \[ a^{[2]} = g\bigl(\vw^{[2]}_1\cdot\va^{[1]} + b^{[2]}_1\bigr) \]
      \textbf{Optional final step:} predict $\yhat=1$ if $a^{[2]}\ge 0.5$, else $\yhat=0$ ---
      exactly as in Week 3.
    \end{column}
  \end{columns}
\end{frame}

\begin{frame}{Notation: superscripts in square brackets}
  \begin{columns}[T,onlytextwidth]
    \begin{column}{0.5\textwidth}
      \begin{defnbox}[Layer notation]
        \small
        \begin{itemize}
          \item $[l]$ --- a quantity belonging to \textbf{layer} $l$.
          \item $(i)$ --- belonging to \textbf{training example} $i$ (as before).
          \item subscript $j$ --- \textbf{unit} $j$ within the layer.
        \end{itemize}
        So $a^{[2](7)}_3$ is the activation of unit 3 of layer 2 on example 7.
      \end{defnbox}
      \footnotesize
      By convention $\va^{[0]}=\vx$, so the input is ``layer 0''. A network is said to have
      $L$ layers when it has $L-1$ hidden layers and one output layer: the input layer is not
      counted.
    \end{column}
    \begin{column}{0.46\textwidth}
      \begin{keyidea}[The one equation of this lecture]
        \[ a^{[l]}_j = g\bigl(\vw^{[l]}_j\cdot\va^{[l-1]} + b^{[l]}_j\bigr) \]
        \footnotesize
        Unit $j$ of layer $l$ takes the \emph{whole} activation vector of the previous layer,
        forms a weighted sum, adds its bias, and applies the activation function $g$.
      \end{keyidea}
      \footnotesize
      For now $g$ is the sigmoid. Week 5 introduces other choices of $g$; in this
      context $g$ is called the \textbf{activation function}.
    \end{column}
  \end{columns}
\end{frame}

\begin{frame}{How many parameters does a layer have?}
  \begin{columns}[T,onlytextwidth]
    \begin{column}{0.5\textwidth}
      \begin{block}{Counting}
        A layer with $n_{\text{in}}$ inputs and $n_{\text{out}}$ units has
        \[ \underbrace{n_{\text{in}}\times n_{\text{out}}}_{\text{weights}}
           + \underbrace{n_{\text{out}}}_{\text{biases}} \]
        parameters: each unit needs one weight per input, plus its own bias.
      \end{block}
      \footnotesize
      \textbf{Demand network}, $4\to3\to1$:\\
      layer 1: $4\cdot3+3=15$; \ layer 2: $3\cdot1+1=4$; \ total $19$.
    \end{column}
    \begin{column}{0.46\textwidth}
      \begin{pitfall}[Count the inputs of \emph{this} layer]
        \small
        Layer 2's units see the $3$ activations of layer 1, not the $4$ original features. A
        frequent exam error is to use $n$ (the number of features) for every layer.
      \end{pitfall}
      \footnotesize
      Parameter count grows with the \emph{product} of adjacent layer widths --- which is why
      a fully connected first layer on a million-pixel image is expensive, and why vision
      networks use other layer types (Week 5).
    \end{column}
  \end{columns}
\end{frame}

% ===============================================================
\section{More Complex Networks and Forward Propagation}

\begin{frame}{A deeper network}
  \begin{columns}[T,onlytextwidth]
    \begin{column}{0.55\textwidth}
      \centering
      \begin{tikzpicture}[x=1.45cm,y=0.62cm]
        \foreach \i in {1,...,4}{\node[inneuron,minimum size=4.5mm] (x\i) at (0,-\i-0.5) {};}
        \foreach \i in {1,...,5}{\node[neuron,minimum size=4.5mm] (a\i) at (1,-\i) {};}
        \foreach \i in {1,...,4}{\node[neuron,minimum size=4.5mm] (b\i) at (2,-\i-0.5) {};}
        \foreach \i in {1,...,3}{\node[neuron,minimum size=4.5mm] (c\i) at (3,-\i-1) {};}
        \node[outneuron,minimum size=4.5mm] (o) at (4,-3) {};
        \foreach \i in {1,...,4}{\foreach \j in {1,...,5}{\draw[synapse,cugray!35] (x\i)--(a\j);}}
        \foreach \i in {1,...,5}{\foreach \j in {1,...,4}{\draw[synapse,cugray!35] (a\i)--(b\j);}}
        \foreach \i in {1,...,4}{\foreach \j in {1,...,3}{\draw[synapse,cugray!35] (b\i)--(c\j);}}
        \foreach \i in {1,...,3}{\draw[synapse,cugray!35] (c\i)--(o);}
        \foreach \x/\l in {0/{$\vx=\va^{[0]}$},1/{$\va^{[1]}$},2/{$\va^{[2]}$},3/{$\va^{[3]}$},4/{$a^{[4]}$}}{
          \node[font=\tiny] at (\x,-6.3) {\l};}
        \foreach \x/\l in {1/layer 1,2/layer 2,3/layer 3,4/layer 4}{
          \node[font=\tiny,cugray] at (\x,-0.1) {\l};}
      \end{tikzpicture}
    \end{column}
    \begin{column}{0.43\textwidth}
      \footnotesize
      \begin{block}{Four layers}
        Three hidden layers and one output layer --- the input is not counted.
      \end{block}
      \begin{block}{Zoom into layer 3, unit 2}
        \[ a^{[3]}_2 = g\bigl(\vw^{[3]}_2\cdot\va^{[2]} + b^{[3]}_2\bigr) \]
        $\vw^{[3]}_2$ has $4$ entries, because layer 2 has $4$ units.
      \end{block}
      \begin{keyidea}
        Every layer takes a vector in and gives a vector out. Stacking layers is composing
        functions:
        $f(\vx)=g^{[4]}\bigl(g^{[3]}\bigl(g^{[2]}\bigl(g^{[1]}(\vx)\bigr)\bigr)\bigr)$.
      \end{keyidea}
    \end{column}
  \end{columns}
\end{frame}

\begin{frame}{Handwritten digits: 0 or 1?}
  \begin{columns}[T,onlytextwidth]
    \begin{column}{0.36\textwidth}
      \centering
      \begin{tikzpicture}[scale=0.3]
        \foreach \i in {0,...,7}{\foreach \j in {0,...,7}{\fill[white] (\i,\j) rectangle ++(1,1);}}
        \foreach \c in {(3,1),(3,2),(3,3),(3,4),(3,5),(3,6),(4,1),(4,2),(4,3),(4,4),(4,5),(4,6),(2,5),(4,7),(3,7)}{
          \fill[cublue] \c rectangle ++(1,1);}
        \draw[cugray!60] (0,0) grid (8,8);
      \end{tikzpicture}\par
      {\scriptsize an $8\times8$ ``1'': $\vx\in\R^{64}$}
      \vspace{2mm}

      \begin{tikzpicture}[node distance=3mm,
          lay/.style={draw=cublue,fill=culight,rounded corners,font=\tiny,
                      minimum height=5mm,inner sep=2pt}]
        \node[lay,draw=cuteal,fill=cuteal!12] (x) {64};
        \node[lay,right=of x] (l1) {25};
        \node[lay,right=of l1] (l2) {15};
        \node[lay,right=of l2,draw=cuamber,fill=cuamber!15] (l3) {1};
        \foreach \a/\b in {x/l1,l1/l2,l2/l3}{\draw[synapse] (\a)--(\b);}
      \end{tikzpicture}
    \end{column}
    \begin{column}{0.6\textwidth}
      \footnotesize
      \begin{block}{Forward propagation}
        \vspace{-2mm}
        \begin{align*}
          \va^{[1]} &= g\bigl(\vw^{[1]}_j\cdot\vx + b^{[1]}_j\bigr)_{j=1..25}
             && \text{25 numbers}\\
          \va^{[2]} &= g\bigl(\vw^{[2]}_j\cdot\va^{[1]} + b^{[2]}_j\bigr)_{j=1..15}
             && \text{15 numbers}\\
          a^{[3]} &= g\bigl(\vw^{[3]}_1\cdot\va^{[2]} + b^{[3]}_1\bigr)
             && \text{1 number}
        \end{align*}
        \vspace{-4mm}

        $f(\vx)=a^{[3]}$; predict ``1'' if $a^{[3]}\ge0.5$.
      \end{block}
      \begin{keyidea}[Inference]
        Computing activations left to right, input to output, is \textbf{forward
        propagation}. It is all that \emph{inference} --- prediction with a trained network
        --- requires. Learning the weights uses \emph{backward} propagation (Week 5).
      \end{keyidea}
      Parameters: $64\cdot25+25=1625$, \ $25\cdot15+15=390$, \ $15+1=16$; \ total $2031$.\par
      Layers usually narrow towards the output.
    \end{column}
  \end{columns}
\end{frame}

\begin{frame}{Forward propagation by hand: a $2\to3\to1$ network}
  \footnotesize
  \begin{columns}[T,onlytextwidth]
    \begin{column}{0.4\textwidth}
      \begin{block}{Parameters}
        Layer 1:\vspace{-1mm}
        \[ \vw^{[1]}_1=\begin{bmatrix}1\\2\end{bmatrix},\
           \vw^{[1]}_2=\begin{bmatrix}-3\\4\end{bmatrix},\
           \vw^{[1]}_3=\begin{bmatrix}5\\-6\end{bmatrix} \]
        \vspace{-2mm}
        \[ b^{[1]}_1=-1,\quad b^{[1]}_2=1,\quad b^{[1]}_3=2 \]
        Layer 2:\vspace{-1mm}
        \[ \vw^{[2]}_1=[\,2,\,-1,\,-3\,],\quad b^{[2]}_1=0.5 \]
      \end{block}
      Input: $\vx=[1,\,2]$.
    \end{column}
    \begin{column}{0.56\textwidth}
      \renewcommand{\arraystretch}{1.2}
      \begin{tabular}{@{}lll@{}}
        \toprule
        unit & $z = \vw\cdot\va_{\text{in}}+b$ & $a=g(z)$\\
        \midrule
        $a^{[1]}_1$ & $1\cdot1+2\cdot2-1=4$       & $0.9820$\\
        $a^{[1]}_2$ & $-3\cdot1+4\cdot2+1=6$      & $0.9975$\\
        $a^{[1]}_3$ & $5\cdot1-6\cdot2+2=-5$      & $0.0067$\\
        \midrule
        $a^{[2]}$   & $2(0.9820)-0.9975-3(0.0067)+0.5$ & \\
                    & $\;=1.9640-0.9975-0.0201+0.5$ & \\
                    & $\;=1.4464$                       & $\mathbf{0.8094}$\\
        \bottomrule
      \end{tabular}
      \vspace{1mm}

      $a^{[2]}=0.8094\ge0.5$, so $\yhat=1$.\par
      \vspace{1mm}
      For $\vx=[-1,-1]$: $\vz^{[1]}=[-4,0,3]$, $\va^{[1]}=[0.0180,0.5000,0.9526]$,
      $z^{[2]}=-2.8217$, $a^{[2]}=0.0562$, so $\yhat=0$.
      \begin{pitfall}[Sanity anchor]
        \footnotesize $g(0)=0.5$ exactly, $g(4)\approx0.982$, $g(-5)\approx0.007$. Check
        hand calculations against these before trusting a long chain.
      \end{pitfall}
    \end{column}
  \end{columns}
\end{frame}

% ===============================================================
\section{Wrap-up and Assessment}

\begin{frame}{Summary}
  \begin{columns}[T,onlytextwidth]
    \begin{column}{0.5\textwidth}
      \begin{block}{What to remember}
        \footnotesize
        \begin{itemize}
          \item A unit computes $a=g(\vw\cdot\va_{\text{in}}+b)$ --- a logistic regression.
          \item A layer is many units sharing the same input vector; it maps a vector to a
                vector.
          \item Hidden layers compute features that training, not the designer, chooses.
          \item Deeper layers build more abstract features from earlier ones.
          \item Forward propagation: compute $\va^{[1]},\va^{[2]},\dots$ left to right.
        \end{itemize}
      \end{block}
    \end{column}
    \begin{column}{0.46\textwidth}
      \begin{block}{Formulas}
        \[ a^{[l]}_j = g\bigl(\vw^{[l]}_j\cdot\va^{[l-1]} + b^{[l]}_j\bigr),\quad \va^{[0]}=\vx \]
        \[ \#\text{params} = n_{\text{in}}\,n_{\text{out}}+n_{\text{out}} \]
      \end{block}
      \begin{block}{Next}
        \footnotesize Part 2: the same computation in TensorFlow, then in NumPy from scratch,
        and finally as a single matrix multiplication.
      \end{block}
    \end{column}
  \end{columns}
\end{frame}

\begin{frame}{Descriptive questions}
  \footnotesize
  \begin{enumerate}
    \item Describe the artificial neuron: its inputs, parameters, computation and output.
          Explain why a single sigmoid unit is identical to logistic regression. \hfill\clo{5}
    \item Give two reasons neural networks became dominant after about 2010, and explain the
          performance-versus-data picture that motivates large networks. \hfill\clo{5}
    \item Using the demand-prediction example, explain the terms input layer, hidden layer,
          output layer and activation. Why is a hidden layer called ``hidden''? \hfill\clo{6}
    \item Contrast feature engineering (Week 2) with feature learning. What does the designer
          still choose when using a neural network? \hfill\clo{6}
    \item Describe what the successive hidden layers of a face-recognition network are
          typically found to compute, and explain why such descriptions are post hoc.
          \hfill\clo{5}
    \item Explain the meaning of each index in $a^{[2](7)}_3$, and state the general
          activation equation for unit $j$ of layer $l$. \hfill\clo{5}
  \end{enumerate}
\end{frame}

\begin{frame}{Descriptive questions (continued)}
  \footnotesize
  \begin{enumerate}\setcounter{enumi}{6}
    \item How many layers does a network with input $\vx$, three hidden layers and one output
          unit have? State the convention you use. \hfill\clo{4}
    \item A network has layer sizes $64\to25\to15\to1$. Compute the number of parameters in
          each layer and in total. Show your working. \hfill\clo{6}
    \item What is forward propagation, and why is it all that is needed for inference? What
          additional computation does training require? \hfill\clo{5}
    \item For the $2\to3\to1$ network on the worked-example slide, compute $a^{[2]}$ for
          $\vx=[0,0]$ and state the prediction. \hfill\clo{8}
    \item Write the output of a three-layer network as a composition of functions, and
          explain why each layer can be viewed as a feature transformation for the next.
          \hfill\clo{6}
    \item In the worked example, $a^{[1]}_2$ is near $1$ for both inputs shown. What does that
          suggest about how useful unit 2 is for distinguishing them? \hfill\clo{5}
  \end{enumerate}
\end{frame}

\begin{frame}{Design questions}
  \footnotesize
  \begin{enumerate}
    \item \textbf{Design a demand model.} A Chattogram garment exporter wants to predict
          whether a new design will sell above 10\,000 units. Propose the input features, an
          architecture (number of layers and units) and the meaning you hope some hidden units
          will acquire. Justify why a network is better here than hand-engineered logistic
          regression --- or argue that it is not. \hfill\clo{12}
    \item \textbf{Design for a parameter budget.} A microcontroller can store at most $5000$
          parameters. Design a network for $28\times28$ greyscale digit images (0 vs.\ 1)
          that fits, and show the count. What do you give up to fit? \hfill\clo{10}
    \item \textbf{Design a check for an inference engine.} You inherit code that claims to do
          forward propagation. Design three hand-computable test cases that would expose,
          respectively: a missing bias, the wrong previous-layer input, and a missing
          activation function. \hfill\clo{12}
    \item \textbf{Design an interpretation study.} A colleague claims hidden unit 7 of a
          trained X-ray classifier ``detects fractures''. Design an experiment to support or
          refute the claim, and state what result would refute it. \hfill\clo{10}
  \end{enumerate}
\end{frame}

\begin{frame}{Reading}
  \begin{block}{Primary text}
    G.~James, D.~Witten, T.~Hastie and R.~Tibshirani,
    \emph{An Introduction to Statistical Learning}, 2nd ed., Springer, 2021.\\[2pt]
    \hfill{\small$\rightarrow$ \textbf{\S10.1--10.2} --- single-layer neural networks, hidden
    units as learned features, and multilayer networks on the MNIST digits.}
  \end{block}
  \begin{block}{Supplementary}
    C.~M.~Bishop, \emph{Pattern Recognition and Machine Learning}, Springer, 2006.\\[2pt]
    \hfill{\small$\rightarrow$ \textbf{\S5.1} --- feed-forward network functions: the
    two-layer network written out unit by unit, and forward propagation.}
  \end{block}
  \vspace{1mm}
  \begin{center}
    \small\color{cugray}
    Lab (CSE 816, Week 4): neurons and layers; coffee roasting in TensorFlow and in NumPy;
    practice lab on neural networks for binary classification.\\
    Practice quizzes: \emph{Neural networks intuition}; \emph{Neural network model}.
  \end{center}
\end{frame}

\end{document}
